% 365_Fundament_FFGFT_En_ch.tex
% Johann Pascher, 15 September 2026
% The foundation of FFGFT in four layers

\providecommand{\BM}{\textbf{[B]}}
\providecommand{\KM}{\textbf{[K]}}
\providecommand{\SM}{\textbf{[S]}}
\providecommand{\SetzM}{\textbf{[POSIT]}}
\providecommand{\QM}{\textbf{[Q]}}
\providecommand{\EM}{\textbf{[E]}}
\providecommand{\XM}{\textbf{[X]}}

\chapter*{The Foundation of FFGFT in Four Layers}
\addcontentsline{toc}{chapter}{The Foundation of FFGFT in Four Layers}

\begin{abstract}
FFGFT rests on four layers of different epistemic character. Layer~1 (the
fundamental equation and the principle of relativity) and Layer~2 (the
electromagnetic sector with $\alpha=1$) are parameter-free in natural units
and generate the structure. Layer~3 ($\xi$ as scaling parameter) translates
the structure into measurable quantities; it is the only free parameter.
Layer~4 (algebraic confirmations: Galois fields, Hodge theory, representation
theory) demonstrates the coherence of the system against known mathematics but
adds no new inputs. Translation into SI units and comparison with the Standard
Model are communication work for established systems, not structural inputs.
The cosmic sector ($H_0/\Lambda$) is bracketed because the Standard Model
derives it equally little.
\end{abstract}

\section*{List of symbols}
\addcontentsline{toc}{section}{List of symbols}

The table explains symbols appearing in more than one section.
Quantities used only locally are introduced there.

\medskip
\textbf{Epistemic markers:}
\begin{tabular}{@{}ll@{}}
\textbf{[B]} & algebraically proven (Python script, exact arithmetic) \\
\textbf{[K]} & numerically confirmed \\
\textbf{[S]} & structurally plausible, identification open \\
\textbf{[POSIT]} & motivated input, not derived \\
\textbf{[E]} & established external mathematics or physics \\
\textbf{[X]} & checked and negative \\
\end{tabular}

\medskip
\begin{tabular}{@{}lp{11cm}@{}}
\toprule
\textbf{Symbol} & \textbf{Meaning} \\
\midrule
\multicolumn{2}{l}{\textit{Fundamental equation}} \\
$\tilde{T}$ & intrinsic time scale of a system \\
$m$ & mass (in natural units = energy) \\
$\tilde{T}\cdot m=1$ & fundamental equation of FFGFT \\
\midrule
\multicolumn{2}{l}{\textit{Geometry}} \\
$T^4$ & four-dimensional torus (compact carrier) \\
$D_4$ & root lattice in $\mathbb{R}^4$; kissing number~24 \\
$\mathbb{Z}_3$ & identification group; triality on $D_4$ \\
$T^4/\mathbb{Z}_3$ & orbifold; nine fixed points \\
\midrule
\multicolumn{2}{l}{\textit{Parameters and fields}} \\
$\xi=4/30\,000$ & dimensionless scaling parameter \\
$\alpha$ & fine-structure constant; in FFGFT units $\alpha=1$ \\
$GF(3^k)$ & finite field of characteristic~3; mode classification \\
$\varphi=(1{+}\sqrt5)/2$ & golden ratio; primitive element of $GF(9)^*$ \\
$\theta=2/9$ & output of the $\mathbb{Z}_3$-circulant diagonalisation; used in §~\ref{sec:schicht1} \\
\midrule
\multicolumn{2}{l}{\textit{Field equations}} \\
$F^{\mu\nu}$ & electromagnetic field strength tensor \\
$j^\nu$ & four-current \\
$\slashed\partial=\gamma^\mu\partial_\mu$ & Feynman slash notation \\
$\hat{H}$ & Hamiltonian \\
\midrule
\multicolumn{2}{l}{\textit{Operations}} \\
Type~I & duality decompactification; information-preserving (limit $R\to\infty$) \\
Type~II & geometric projection; lossy (e.g.\ projection onto subspaces) \\
Type~III & representational translation; bijective \\
\bottomrule
\end{tabular}


% ============================================================
\section*{On the intention}
% ============================================================

FFGFT does not make predictions in the sense of asserting something
contentually new. It rewrites known relations so that hidden structure becomes
visible. Experimental confirmation is not the goal --- the reformulated equations
are identical to the established ones and therefore trivially correct. What is
interesting is what the rewriting reveals: that $H_0$ follows from the geometry,
that $a_0=cH_0/(2\pi)$ is not a coincidence, that $\theta=2/9$ was not an input.
Had the cosmic sector been unknown, the geometry would have supplied it --- one
could then measure whether it agrees. That it is already known does not make the
statement false, but it makes testability harder, because one can never be fully
certain whether an exponent was found or chosen.

Confirmations do not succeed --- that is not a deficiency but the consequence of
the method. Whoever seeks confirmation has misunderstood the intention.

The gain over the Standard Model lies not in new numerical values but in the
reduction of free parameters: the Standard Model uses $\sim$19 independent
constants; FFGFT manages with \emph{one} dimensionless parameter $\xi$, which
itself follows from Galois group orders. What the Standard Model treats as
independent measured quantities --- $\alpha$, mass ratios, mixing angles ---
appears in FFGFT as a consequence of the torus geometry and field structure.
That is the structural gain: not a new prediction, but a deeper explanation
of existing ratios.

% ============================================================
\section{Layer 1: the fundamental equation}
% ============================================================

In natural units ($\hbar=c=1$) the only fundamental equation of FFGFT is
\[
  \tilde{T} \cdot m = 1 .
\]
It states: time scale and mass are reciprocal representations of the same object.
Combined with the principle of relativity (Lorentz invariance, causal structure,
action principle) this implies:
\begin{itemize}[nosep]
\item A compact carrier $T^4$ as the smallest structure compatible with
      $\tilde{T}\cdot m=1$, Lorentz invariance and three spatial dimensions \SetzM.
\item $\mathbb{Z}_3$ as the most economical identification with three generations \SetzM.
\item $D_4$ as the densest lattice in four dimensions with $\mathbb{Z}_3$ triality \SetzM.
\end{itemize}
These three posits are motivations, not derivations; they are vindicated by
internal coherence (Dok.~248).

\subsection*{Three levels of description}

It is important to distinguish what belongs to the geometry of the torus and
what is a translation into other mathematical languages (Dok.~270/330):

\begin{enumerate}[nosep]
\item \textbf{Torus geometry (Level~1)} --- the physics itself. $T^4/\mathbb{Z}_3$
on the $D_4$ lattice, modes as winding numbers, fixed points, ratios. Static,
in $\mathbb{Q}$ and algebraic fields.
\item \textbf{Galois fields (Level~2)} --- the algebraic classification of the
modes. $GF(9)$, $GF(27)$, $GF(729)$ describe which mode orders can exist; the
GF tower is the algebraic language of Level~1, not a separate level of reality.
Bijective translation (Type~III, Dok.~270): no information loss, no new content.
\item \textbf{Matrix/Hilbert space (Level~3)} --- the QM language.
$T^4\leftrightarrow\mathcal{H}$ via Fourier on the torus; likewise Type~III,
bijective. Type~I denotes duality decompactifications ($R\to\infty$);
Type~II geometric projections with information loss. Enables connection to Schrödinger, Dirac, product states --- but adds
no physics. The Lagrangian and field equations belong to this level as imported
vocabulary \EM\ (R120), to keep the theory compatible with established frameworks.
\end{enumerate}

All three levels describe the same static snapshot in different languages.
What changes is the language, not the content.

\subsection*{Static picture --- snapshot}

The theory describes a static structure: a snapshot. Time is a coordinate of
the torus $T^4$, not a distinguished evolution parameter; $\tilde{T}\cdot m=1$
is a property of each state, not a propagation law. What appears as a temporal
evolution is the transition between snapshots --- which the theory does not
describe in Layers~1--3. The position-dependent time field
$T(x,t)\equiv1/m(x,t)$ is the snapshot relation at point $x$: when the mass
$m(x,t)$ varies in a gravitational potential, $T(x,t)$ varies with it
necessarily --- not because $T$ is an independent dynamical field, but because
$T=1/m$ holds at every instant. The ``extended'' Schrödinger equation with
$T(x,t)$ formulated in the early documents is therefore a rewriting of the
mass-position-dependent coupling to the gravitational potential, not a new
degree of freedom.

\subsection*{Ratio-based computation and exact rational arithmetic}

Since $\tilde{T}\cdot m=1$ in natural units formulates all physical statements as
ratios --- mass to mass, time to time, length to length --- the fundamental
calculations remain in $\mathbb{Q}$ or finite algebraic extensions such as
$\mathbb{Q}(\sqrt5)$. No floating-point operations are required as long as no SI
anchor is inserted. The FFGFT verification scripts therefore use exact rational
arithmetic (SymPy \texttt{Rational}, Python's \texttt{Fraction}); replay
identity --- bit-identical reproducibility of every step --- is the natural
test criterion. Rounded approximations arise only when translating into SI units
(Layer~3).

\paragraph{What already follows.}
From $\tilde{T}\cdot m=1$ alone follow time dilation as a mass effect, redshift
as path lengthening in the intrinsic time field and the basic structure of
gravitation (Dok.~078/270). The pattern of the mass ladder follows from the
geometry of $T^4$; $\theta=2/9$ is the output of the $\mathbb{Z}_3$-circulant
diagonalisation, not an input (Dok.~230/231, \enquote{found, not sought}).

\paragraph{Gravity as a derived quantity.}
Gravitational effects follow from $\tilde T\cdot m=1$ applied locally: if the
mass $m(x)$ varies in a gravitational field, $\tilde T(x)=1/m(x)$ varies with
it necessarily. There is no independent gravitational field and no graviton as
an independent degree of freedom; gravity is a consequence of the
position-dependent time field, not a separate foundation (Dok.~078 \KM).
The choice between ``mass constant, space curved'' (Einstein) and ``space
flat, mass variable'' (FFGFT) is a \emph{choice of representation} ---
exactly equivalent and interconvertible by a coordinate transformation,
like the Schrödinger and Heisenberg pictures in QM. The physically most
natural reading within $\tilde T\cdot m=1$ is that \emph{both} vary and
that holding either one constant is a simplifying assumption (Dok.~078).

\paragraph{No gap between relativity and quantum theory.}
The apparent tension between general relativity and quantum mechanics is an
artefact of the level of description. At Level~1 (torus geometry) this
opposition does not exist: GR and QM are two Type~III translations of the
same static snapshot --- GR into the language of spacetime curvature, QM
into the language of operators on $\mathcal{H}$. FFGFT is not a new theory
that \emph{unifies} them; it shows that the opposition does not exist at
Level~1.

FFGFT achieves more than standard QM describes. In standard QM the time $t$
is an external parameter outside the system --- the wave function evolves in
a prescribed background time that is itself not quantised. In FFGFT time
belongs to the system: $\tilde T=1/m$ is a property of the physical object,
not a background structure. This changes the status of quantum randomness:
what appears in standard QM as fundamental, ontological chance is in FFGFT
deterministic time-field dynamics --- in principle determined, in practice
unpredictable (Dok.~020 \KM). There is no collapse: wave functions evolve
continuously; the Born rule emerges from the time-field interaction instead
of being posited as an axiom. Randomness is epistemic, not ontological.

% ============================================================
\section{Layer 2: the electromagnetic sector and simplification of the field equations}
% ============================================================

\subsection*{FFGFT units (formerly: T0 units)}

Setting $\alpha=1$ in natural units is standard practice in theoretical physics
(Heaviside-Lorentz units, Gaussian units) and not a FFGFT-specific decision.
The FFGFT unit system additionally sets $c=\hbar=\alpha=G=1$ (Dok.~011;
formerly called T0 units). $\alpha=1$ is not a posit in the sense of R56 but a
\emph{choice of units}: the electromagnetic coupling strength remains the same
physical quantity; it receives the numerical value~1. The SI value $1/137$ does
not disappear --- it is fully reconstructible from $\xi$ and $E_0$:
$\alpha_{\mathrm{SI}}=\xi\cdot(E_0/1\,\mathrm{MeV})^2\approx1/137$ (Dok.~011
\KM). In FFGFT units and in SI units $\alpha$ is the same quantity --- only its
numerical representation changes.

\subsection*{Derived electromagnetic quantities}

With $\alpha=1$ electromagnetic quantities become derived forms of mass and time:
\begin{align*}
  e &= 1 \quad \text{(charge dimensionless)},\\
  V &= E/e = E \quad \text{(voltage = energy)},\\
  I &= e/T = 1/T = m \quad \text{(current = mass)}.
\end{align*}
The Coulomb potential simplifies to $V(r)=-1/r$.

\subsection*{Maxwell, Schrödinger and Dirac equations in FFGFT units}

The established field equations take their most geometrically transparent form
in FFGFT units --- not because they are changed, but because the unit choice
sets the coupling constants to $1$ and thereby makes visible the structure that
was always there:
\begin{align*}
\text{Maxwell:}\quad &
  \partial_\mu F^{\mu\nu} = j^\nu, \qquad
  \partial_{[\mu}F_{\nu\rho]}=0,\\
\text{Schrödinger:}\quad &
  i\,\partial_t\,\psi = \left(-\tfrac12\nabla^2 + V\right)\psi\\
\text{Dirac:}\quad &
  (i\slashed\partial - m)\,\psi = 0.
\end{align*}
$\hbar$, $c$, $e$ and $4\pi\varepsilon_0$ no longer appear explicitly; they
are absorbed into the unit choice. This is not new physics but the known physics
in a notation that makes the mass structure directly readable: $i\,\partial_t\,\psi=m\,\psi$
in the free case, $\partial_\mu F^{\mu\nu}=j^\nu$ without a prefactor.

\paragraph{What becomes concrete.}
Only with Layer~2 does the theory become concrete enough for particle physics:
charge quantisation \BM, $\mathrm{su}(3)$ from $\mathbb{Z}_3$ triality \BM,
charge definition and gauge structure. The mass layers of the lepton sector are
derivable from $\xi$ \KM; the mass ratio $m_\mu/m_e$ serves as a declared
reference point \SetzM\ (Dok.~190 P42) to keep the layer independently checkable.

% ============================================================
\section{Layer 3: the only free parameter}
% ============================================================

$\xi = 4/30\,000$ appears as the only dimensionless parameter but is not free
given the Galois group orders of the GF tower. From the winding numbers of the
$T^4$ torus (derived from char$(\text{GF}(3^n))$, $|\text{GF}(3)^*|$ and the
minimal prime divisor of $|\text{GF}(81)^*|$) and the field orders
$|\text{GF}(9)^*|=8$, $|\text{GF}(27)|=27$, $5=\text{min.prime}(80)$, $\xi$
follows algebraically \KM:
\[
  \xi = \frac{(r_\mu/r_e)^2}{8^2\cdot5^2\cdot27} = \frac{(12/5)^2}{43200} = \frac{4}{30\,000}.
\]
(Dok.~336/338; verification scripts pruef\_330, pruef\_331.) The premises of
this derivation --- the winding-number recursion and $m_\mu/m_e$ as reference
point (Dok.~190 P42) --- carry a $[\text{K}]$ status, not $[\text{B}]$. $\xi$
is therefore not completely free, but not yet completely derived.
Note: the value $m_\mu/m_e = 12/5 = 2.4$ used in Dok.~338 is the
\emph{structural ratio} of the winding numbers (Level~1), not the measured
mass ratio $\approx 206.77$; the measured ratio follows from the structural
value after scaling by $\xi$ and the Galois-field exponents.

With $\xi$, Layers~1 and~2 yield measurable scales: Compton wavelength,
Hubble length as geometric maximal scale (not an expansion horizon),
$a_0=cH_0/(2\pi)$ \KM.

\paragraph{What is bracketed.}
The cosmic sector ($H_0$, $\Lambda$) is bracketed. The Standard Model derives
$H_0$ equally little; the comparison would be asymmetric (Dok.~190 P39). The
exponent~10 in $H_0\propto\xi^{10}$ is a posit (Dok.~190 P20); it is
structurally motivated but not forced.

\paragraph{SI anchors are comparison tools.}
$m_e$, $\bar\lambda_e$, $\nu_{\mathrm{Cs}}$ and all other SI anchors serve
exclusively for translation into known unit systems. Every dimensionless
statement of FFGFT is a function of $\xi$ and the posits; no SM value is a
structural input. The complexity of comparison tables and precision calculations
is communicative, not structural --- a concession to established systems that
cannot be abandoned without notice.

% ============================================================
\section{Layer 4: algebraic confirmations}
% ============================================================

The Galois fields $GF(9)$, $GF(27)$, $GF(729)$, Hodge theory on
$T^4/\mathbb{Z}_3$, representation theory of $\mathrm{Aut}(D_4)$ and
comparisons with external carriers (Dok.~361--364) are \emph{consistency
checks}, not new inputs. They show that the structure generated by
$\tilde{T}\cdot m=1$ and the three posits is coherent against known
mathematics; they could all have followed from $\xi$ and the posits alone
and in a sense do.

The icosahedron (Dok.~364) is not a special case: every geometric body with
non-crystallographic symmetry encounters the same structure --- lattice-compatible
part in $\mathrm{Aut}(D_4)$, remainder as phase in $GF(3^k)$. This is not
coincidence but an expression of the prime~3 as fundamental building block:
what cannot exist as a root of unity in characteristic~3 must be a lattice
operation, and vice versa.

\section*{Tool and method}
\addcontentsline{toc}{section}{Tool and method}

Claude (Anthropic) was used as a tool for: writing and executing the Python
verification scripts; drafting the \LaTeX\ sources from results and corrections
of the author; maintaining the changelog and cross-references. The tool computes
and writes; the author decides what is computed and written, judges every result
and determines what it means physically. Every algebraic statement is backed by
a Python verification script with explicit assertions --- if an assertion fails,
the statement is not used.

% ============================================================
\section{Summary: what FFGFT needs and what it does not}
% ============================================================

\begin{center}
\begin{tabular}{@{}lll@{}}
\toprule
Layer & Content & Status \\
\midrule
1 & $\tilde{T}\cdot m=1$ + principle of relativity & fundamental law + 3 posits \\
2 & $\alpha=1$, charge definition & posit, EM sector follows \\
3 & $\xi=4/30\,000$ & only free parameter \\
4 & Galois, Hodge, representation theory & consistency confirmations \\
\midrule
--- & SI anchors, SM comparisons & translation, not input \\
--- & cosmic sector ($H_0$, $\Lambda$) & bracketed (SM likewise) \\
\bottomrule
\end{tabular}
\end{center}

What FFGFT does \emph{not} need: the Standard Model, its constants, its
Lagrangian as input (only for comparison), external calibrations beyond $\xi$.
What it delivers: a self-consistent structure that is parameter-lean in natural
units and reproduces all known ratios in SI units. It becomes testable only when
anchors are inserted --- but the anchors are translation, not foundation.

\begin{thebibliography}{9}

\bibitem{Planck1900}
M.~Planck,
\textit{\"Uber irreversible Strahlungsvorg\"ange},
Sitzungsber.\ Preu\ss.\ Akad.\ Wiss.\ 5 (1900) 440--480.
Natural units ($\hbar=c=G=1$) as a concept.

\bibitem{Heaviside1893}
O.~Heaviside,
\textit{Electromagnetic Theory}, vol.~1,
The Electrician, London, 1893.
Rationalised units, $\alpha=1$ convention (Heaviside-Lorentz).

\bibitem{Lidl1997}
R.~Lidl, H.~Niederreiter,
\textit{Finite Fields},
2nd ed., Cambridge University Press, 1997.
$\xi$ as a Galois number from group orders of $GF(9)^*$, $GF(27)$.

\bibitem{Pascher011}
J.~Pascher,
\textit{Dok.~011: Fine-structure constant in FFGFT units},
FFGFT corpus.

\bibitem{Pascher019}
J.~Pascher,
\textit{Dok.~019: Lagrangian formulation},
FFGFT corpus.

\bibitem{Pascher270}
J.~Pascher,
\textit{Dok.~270: Projection chain $T^4\to T^0$},
FFGFT corpus.

\bibitem{Pascher330}
J.~Pascher,
\textit{Dok.~330: Operations on $T^4$, demarcations},
FFGFT corpus.

\bibitem{Pascher336}
J.~Pascher,
\textit{Dok.~336: GF(9)-FFGFT bridge},
FFGFT corpus.

\bibitem{Pascher338}
J.~Pascher,
\textit{Dok.~338: Galois masses and FFGFT},
FFGFT corpus.

\end{thebibliography}
